% !TEX root = ../main.tex
\chapter{Counterparty Risk, Margin and CVA}\label{ch:counterparty}
\begin{paracol}{2}

\begin{sources}
\begin{tabularx}{\linewidth}{@{}l l L@{}}
\toprule
\textbf{Lecture} & \textbf{Speaker} & \textbf{Content}\\
\midrule
2024 L10 & J.~Shepherd & \emph{Counterparty Risk Optimization} (Quantile Technologies, part of LSEG).
Types of risk and central clearing; coherent risk measures, VaR and expected shortfall (historical
and model-building estimates, time scaling, back-testing, Rockafellar--Uryasev); variation and initial
margin, ISDA SIMM; multilateral optimisation of initial margin in a network of banks, numerical issues,
``nice'' and ``fair'' solutions.\\
2013 L26 & Y.~Tang & \emph{Introduction to Counterparty Credit Risk} (Morgan Stanley, given remotely).
CVA, DVA, funding risk, wrong-way risk (Berkshire Hathaway's index puts), the CVA conundrum;
enterprise-level derivatives modelling; martingale measures, martingale testing, resampling and
interpolation. Last minutes: closing remarks of the instructors (\cref{ch08:sec-closing}).\\
\bottomrule
\end{tabularx}

\smallskip
\textbf{Files.} Slides \file{mit18_642_f24_lec10.pdf} (51 slides); 2013 slides
\file{MIT18_S096F13_lecnote26.pdf} (27 slides).
\textbf{Problem sets.} No problem set of either year covers this chapter; all exercises are marked
``not from the course''.
\textbf{Not covered in class.} 2024: slide 17 (FRTB) was skipped, slides 26--28 (SIMM practicalities)
and 48 (observed savings) only shown; 2013: slides 21--25 (extinguisher trade, ``bigger picture'')
were not reached for lack of time (summarised in \cref{ch08:rem-extinguisher}).
\end{sources}

\section*{Motivation}
An over-the-counter (OTC) derivative is a bilateral promise. Over the life of a swap sometimes you owe
your counterparty money and sometimes it owes you; if it defaults while it owes you, you lose part of the
receivable. Before 2008 this \emph{counterparty credit risk} was a side issue. During the crisis the
mark-to-market losses it caused were larger than the losses from actual defaults \ts{13}{26}{6:22}, and the
regulatory response (Basel~III capital for CVA, mandatory central clearing, margin rules for uncleared
derivatives) changed the whole derivatives business. The two guest lectures show the two sides of this change:
\begin{itemize}
  \item \textbf{2013 (Y.~Tang, Morgan Stanley)}: \emph{pricing} counterparty risk. The credit valuation
        adjustment (CVA) and its relatives (DVA, funding), wrong-way risk, and why all of them need models of
        the whole portfolio, not of one trade at a time.
  \item \textbf{2024 (J.~Shepherd, Quantile/LSEG)}: \emph{mitigating} it with collateral. Variation margin,
        initial margin computed as a value-at-risk, and how a network of banks can reduce the total initial
        margin they post by agreeing on new trades, found by convex optimisation.
\end{itemize}
The mathematics: option-like expectations under the pricing measure, coherent risk measures (VaR and
expected shortfall), quadratic forms of Gaussian vectors, convex and mixed-integer optimisation, and change of
numeraire. VaR is treated in more depth, from the 2013 lecture on the topic, in \cref{ch:var}.

% ==================================================================
\section{Counterparty risk in OTC derivatives}\label{ch08:sec-ccr}

\paragraph{Types of risk.} Trading derivatives creates several risks \ts{24}{10}{3:11}:
\emph{market} risk (values move with the market), \emph{credit} risk (a debtor does not pay an obligation),
\emph{operational and legal} risk (people, systems and events; the hardest to model, because it is mostly
people making mistakes), \emph{liquidity} risk (a transaction cannot be executed at the market price, typically
because the position is so large that selling it moves the price, or outgoing cash flows cannot be funded),
and \emph{counterparty} risk: the potential loss on the market value of transactions when the counterparty
defaults.

\paragraph{Counterparty versus credit risk.} Suppose you buy protection (a credit default swap, CDS) on a
JPMorgan bond from Goldman Sachs. You have \emph{credit} risk to JPMorgan: the value of the CDS depends on
JPMorgan defaulting. You have \emph{counterparty} risk to Goldman Sachs: whether it will actually pay what the
CDS is worth \ts{24}{10}{4:13}. For a corporate bond the economics are similar, but the risk to the issuer
is called \emph{issuer} risk \ts{13}{26}{3:10}. Counterparty risk has two components \ts{13}{26}{2:04}:
\begin{itemize}
  \item \emph{default risk}: at default we recover only a fraction $R$ of what we are owed;
  \item \emph{mark-to-market (MTM) risk}: even if nothing is due for ten years, a widening of the
        counterparty's credit spread lowers the value of the trade today, because a buyer of the trade would
        pay less for it \ts{13}{26}{3:10}.
\end{itemize}
What makes counterparty risk harder than bond credit risk is that the amount at risk is \emph{random and
two-sided}: the value of a swap can be positive or negative, and it depends on all trades with the same
counterparty.

\paragraph{Mitigation.} The standard tools (slide~4 of \file{lec10}) are \ts{24}{10}{5:15}: \emph{netting}
of offsetting payments; \emph{hedging} with new trades; \emph{collateralisation} (collecting margin, a kind of
insurance payment against the counterparty's default); contractual clauses such as periodic resets of the
mark-to-market value or break clauses; and \emph{central clearing counterparties} (CCPs).

\paragraph{Central clearing.} A CCP steps into the middle of bilateral trades by \emph{novation}: the trade
between A and B is replaced by two trades, A with the CCP and the CCP with B \ts{24}{10}{6:15}. By itself
this only changes the counterparty; the gain comes from \emph{multilateral netting}, since all of a firm's
positions now face the same counterparty and offset (\cref{ch08:fig-ccp}) \ts{24}{10}{7:20}. As a rough
rule, total counterparty risk scales like $\sum_{\text{counterparties}}\|\text{position}\|$, and netting reduces
the number of terms. Standard (vanilla) interest-rate swaps must be cleared by regulation, while more exotic
products such as swaptions are not eligible, so every bank ends up with a mixture of cleared and bilateral
trades \ts{24}{10}{8:25}.

\begin{figure}[ht]
\centering
\begin{tikzpicture}[>=Stealth,every node/.style={font=\small},
  pt/.style={circle,draw,thick,minimum size=7mm,inner sep=0pt}]
\begin{scope}
  \node[pt] (C) at (0,2.4) {C}; \node[pt] (D) at (2.4,2.4) {D};
  \node[pt] (A) at (0,0) {A};   \node[pt] (B) at (2.4,0) {B};
  \draw[->,thick,defcol] (C) -- node[left]{125} (A);
  \draw[->,thick,defcol] (A) -- node[below]{100} (B);
  \draw[->,thick,defcol] (B) -- node[right]{100} (D);
  \draw[->,thick,defcol] (D) -- node[above]{50} (C);
  \node at (1.2,-1.0) {bilateral: gross $375$};
\end{scope}
\begin{scope}[xshift=5.2cm]
  \node[pt] (C) at (0,2.4) {C}; \node[pt] (D) at (2.4,2.4) {D};
  \node[pt] (A) at (0,0) {A};   \node[pt] (B) at (2.4,0) {B};
  \node[circle,draw,thick,fill=okcol!15,minimum size=9mm,inner sep=0pt] (P) at (1.2,1.2) {\scriptsize CCP};
  \draw[thick,thmcol] (A) -- node[pos=0.45,below right,inner sep=1pt]{$+25$} (P);
  \draw[thick,black!30] (B) -- node[pos=0.45,below left,inner sep=1pt]{$0$} (P);
  \draw[thick,thmcol] (C) -- node[pos=0.45,above right,inner sep=1pt]{$-75$} (P);
  \draw[thick,thmcol] (D) -- node[pos=0.45,above left,inner sep=1pt]{$+50$} (P);
  \node at (1.2,-1.0) {after netting at the CCP: $150$};
\end{scope}
\end{tikzpicture}
\caption{Multilateral netting at a CCP (slide~5 of \file{lec10}). Left: four bilateral positions; an arrow
into a node counts $+$, out of a node $-$. Right: after novation each party faces only the CCP and its
positions net: A has $+125-100=+25$, B has $+100-100=0$ (no counterparty risk left), C has $+50-125=-75$,
D has $+100-50=+50$. The CCP itself is flat.}\label{ch08:fig-ccp}
\end{figure}

\begin{finance}[What a CCP costs (not discussed in class)]
Clearing concentrates risk in the CCP, which protects itself by collecting initial margin and contributions to a
default fund from its members. Splitting a portfolio between cleared and bilateral trades also destroys netting
\emph{between} the two sets. This is why ``move more trades into the CCP'' is a genuine design choice in the
optimisation of \cref{ch08:sec-opt}, not an automatic improvement.
\end{finance}

% ==================================================================
\section{Exposure, netting and collateral}\label{ch08:sec-exposure}

Fix a counterparty and a \emph{netting set} (all trades covered by one master agreement, whose values are
offset at default). Let $V(t)$ be the value \emph{to us} of the netting set computed with a
counterparty-default-free model, and $C(t)$ the value of the collateral we hold ($C<0$ if we have posted).

\begin{definition}[Exposure profiles]\label{ch08:def-exposure}
The \emph{exposure} at time $t$ is $E(t)=\bigl(V(t)-C(t)\bigr)^+$, the amount we would lose (before
recovery) if the counterparty defaulted at $t$. Then:
\begin{itemize}
  \item \emph{expected exposure} $\mathrm{EE}(t)=\E[E(t)]$, and \emph{expected positive exposure}
        $\mathrm{EPE}=\frac1T\int_0^T\mathrm{EE}(t)\dd t$, called \emph{mean positive exposure} (MPE) in the 2013
        slides;
  \item \emph{negative exposure} $\bigl(V(t)-C(t)\bigr)^-=\max(C(t)-V(t),0)$, what we owe; its mean is the MNE of
        the 2013 slides;
  \item \emph{potential future exposure} $\mathrm{PFE}_q(t)$, the $q$-quantile of $E(t)$ (typically $q=95\%$
        or $99\%$), used by banks to set credit limits (not discussed in class).
\end{itemize}
For pricing (CVA) the expectation is under the pricing measure $\Q$; for risk management under $\Prob$.
\end{definition}

The positive part is the key non-linearity: if $V>0$ at default we recover only $R\,V$, but if $V<0$ we
still owe the full amount to the defaulted counterparty. There is no ``default gain'', and this asymmetry
makes the loss an \emph{option-like} payoff \ts{13}{26}{19:10}. \Cref{ch08:fig-exposure} shows the typical
shapes: the uncertainty of $V(t)$ grows like $\sqrt t$ by diffusion, while for a swap the remaining cash flows,
and with them the sensitivity to rates, shrink linearly to maturity.

\begin{figure}[ht]
\centering
\begin{tikzpicture}
\begin{axis}[width=0.78\linewidth,height=5.2cm,xmin=0,xmax=10,ymin=0,ymax=2.8,
  xlabel={time $t$ (years)},ylabel={exposure / $\sigma$},
  tick label style={font=\scriptsize},label style={font=\small},
  legend style={font=\scriptsize,at={(0.02,0.98)},anchor=north west,draw=none,fill=none},
  legend cell align=left]
\addplot[defcol,thick,domain=0:10,samples=101]{0.3989*sqrt(x)*(1-x/10)};
\addlegendentry{swap, EE}
\addplot[defcol,thick,dashed,domain=0:10,samples=101]{1.645*sqrt(x)*(1-x/10)};
\addlegendentry{swap, PFE$_{95\%}$}
\addplot[thmcol,thick,domain=0:10,samples=101]{0.3989*0.5*sqrt(x)};
\addlegendentry{forward, EE}
\addplot[thmcol,thick,dashed,domain=0:10,samples=101]{1.645*0.5*sqrt(x)};
\addlegendentry{forward, PFE$_{95\%}$}
\end{axis}
\end{tikzpicture}
\caption{Stylised exposure profiles (not from the course). If $V(t)\sim N(0,s(t)^2)$ and there is no
collateral, then $\mathrm{EE}(t)=s(t)/\sqrt{2\pi}\approx0.40\,s(t)$ and $\mathrm{PFE}_{95\%}(t)=1.645\,s(t)$.
Swap: $s(t)=\sqrt t\,(1-t/T)$ with $T=10$ (diffusion times remaining duration, peak at $T/3$).
FX forward: $s(t)=0.5\sqrt t$ (a single exchange at maturity).}\label{ch08:fig-exposure}
\end{figure}

\begin{proposition}[Netting reduces exposure]\label{ch08:prop-netting}
For trades with values $V_1,\dots,V_n$ in one netting set,
$\bigl(\sum_iV_i\bigr)^+\le\sum_iV_i^+$ pathwise; hence $\mathrm{EE}_{\text{net}}(t)\le\sum_i\mathrm{EE}_i(t)$.
\end{proposition}
\begin{proof}
$x\mapsto x^+$ is subadditive: $(x+y)^+\le x^++y^+$ because the right side is $\ge x+y$ and $\ge0$.
Induct on $n$ and take expectations.
\end{proof}
If one trade goes from $0$ to $+\$100$M and an offsetting trade with the same counterparty goes to $-\$100$M, the
netted default loss is zero \ts{13}{26}{18:05}. The price of counterparty risk is therefore \emph{not} a sum of
per-trade quantities: to price it for one new trade you must know every other trade facing the same
counterparty. Ordinary derivative pricing never needs this \ts{13}{26}{7:29}.

\paragraph{Collateral.} Under a \emph{Credit Support Annex} (CSA) to the master agreement the parties exchange
collateral. \emph{Variation margin} (VM) is exchanged daily and equals the current value, $\mathrm{VM}(t)=V(t)$;
the CSA also fixes thresholds below which no collateral is called, minimum transfer amounts and eligible assets
(these details were not discussed in class, apart from the thresholds mentioned on the SIMM slides). With perfect
daily VM the exposure is no longer the \emph{level} of $V$ but its \emph{increment} between the last margin call
the defaulter honoured and the moment we have closed out or replaced the trades. That increment is covered by
\emph{initial margin} (\cref{ch08:sec-margin}).

\begin{intuition}[From $\sqrt{T}$ to $\sqrt{\text{MPOR}}$]
For a diffusive value process the typical size of $V(t)$ after 5 years is $\sigma\sqrt{5}$, while the typical move
over a 10-business-day close-out period is $\sigma\sqrt{10/250}$, about $9\%$ of it. Collateralisation replaces
a large, slowly growing exposure by a small, constantly renewed one. The price is liquidity: margin must be
paid in cash or high-quality securities, often at short notice.
\end{intuition}

\begin{exercise}[Netting and correlation (not from the course)]
Two trades with the same counterparty have values $V_1,V_2\sim N(0,\sigma^2)$ at time $t$, with correlation $\rho$, and
there is no collateral.
\begin{enumerate}[(a)]
  \item Compute $\mathrm{EE}$ with and without netting and show that the ratio is
      $\sqrt{(1+\rho)/2}$.
  \item What happens for $\rho=-1$ and $\rho=1$?
  \item Explain why the CVA of a new trade depends on the
      existing portfolio (the ``incremental CVA'').
\end{enumerate}
\end{exercise}

% ==================================================================
\section{CVA, DVA and funding: pricing counterparty risk}\label{ch08:sec-cva}

\subsection{A puzzle}
Tang opened with a question \ts{13}{26}{8:32}. You have an uncollateralised swap with value $0$ on day one
(ignoring counterparty risk). Later its value is $+\$100$M, and then the counterparty unexpectedly defaults with
$50\%$ recovery, so you receive $\$50$M in cash. Over the life of the trade, have you made $\$50$M, lost $\$50$M,
or are you flat?

A dealer has \emph{lost} $\$50$M \ts{13}{26}{12:45}. Its market risk is hedged, for example by an offsetting swap
with another dealer or an exchange (or the trade is simply intermediated back to back). When the client trade went
from $0$ to $+100$, the hedge went from $0$ to $-100$, and the hedge counterparty is solvent and must be paid in
full. The dealer receives $50$ and pays $100$. In general the loss at default is $(1-R)$ times the exposure. A
\emph{CVA desk} that had sold the trading desk protection on the counterparty (and itself bought CDS protection)
would have made the trading desk whole \ts{13}{26}{13:53}.

\subsection{Credit valuation adjustment}
\begin{definition}[CVA, 2013 slides]\label{ch08:def-cva}
The \emph{credit valuation adjustment} is the price of counterparty credit risk, an adjustment to the value given
by a counterparty-default-free model or a broker quote \ts{13}{26}{4:13}. With $\tau$ the counterparty's default
time, $T$ the final maturity of the netting set, $R_\tau$ the recovery rate, $V_{\tau-},C_{\tau-}$ the portfolio
value and collateral just before default, and $\beta_t=\exp\bigl(\int_0^tr_u\dd u\bigr)$ the money-market account,
\begin{equation}\label{ch08:eq-cva}
  \mathrm{CVA}_{\mathrm{MPE}} = -\E_0^{\Q}\!\left[\ind_{\{\tau\le T\}}\,
      \frac{(V_{\tau-}-C_{\tau-})^+\,(1-R_\tau)}{\beta_\tau}\right] \ts{13}{26}{15:57}.
\end{equation}
The adjusted value of the portfolio is $V_0+\mathrm{CVA}_{\mathrm{MPE}}$.
\end{definition}
The minus sign matters in practice: CVA on receivables is a cost, the amount charged to the counterparty.
Theoretical papers often define CVA as a positive number and subtract it; that is fine as long as you are
consistent \ts{13}{26}{14:55}. In \eqref{ch08:eq-cva}, $\ind_{\{\tau\le T\}}$ says there is no counterparty risk
if the counterparty defaults after the portfolio has matured; $(V-C)^+$ is the exposure; $(1-R)$ turns it into a
loss given default; and $\beta_\tau$ discounts \ts{13}{26}{17:01}.

\begin{proposition}[CVA as an integral of expected exposure (not derived in class)]\label{ch08:prop-cvaee}
Assume $\tau$ is independent of the market variables (no wrong-way risk, \cref{ch08:sec-wwr}), has a
deterministic hazard rate $h(t)$ with survival function $S(t)=\exp\bigl(-\int_0^th\bigr)$, and $R$ is constant.
Then
\[
  \mathrm{CVA} = -(1-R)\int_0^T \mathrm{EE}^*(t)\,h(t)S(t)\dd t,\qquad
  \mathrm{EE}^*(t)=\E^{\Q}\!\left[\frac{(V_t-C_t)^+}{\beta_t}\right].
\]
If moreover $h$ and $r$ are constant, $\mathrm{EE}(t)\approx\mathrm{EPE}$ is roughly flat, and the CDS spread obeys
the ``credit triangle'' $s\approx(1-R)h$, then
\begin{equation}\label{ch08:eq-cvabote}
  \mathrm{CVA}\approx-\,\mathrm{EPE}\times s\times D,\qquad D=\int_0^Te^{-(r+h)t}\dd t=\frac{1-e^{-(r+h)T}}{r+h},
\end{equation}
where $D$ is the risky duration (annuity) of the counterparty. This is the ``back-of-the-envelope'' formula of
the slides: $-\,$MPE $\times$ CDS par spread $\times$ duration \ts{13}{26}{13:53}.
\end{proposition}
\begin{proof}
Let $g(t)=\E^{\Q}\bigl[(V_t-C_t)^+/\beta_t\bigr]$. Since $\tau$ is independent of the market,
$\E[\ind_{\{\tau\le T\}}(1-R)(V_{\tau-}-C_{\tau-})^+/\beta_\tau]=(1-R)\E[\ind_{\{\tau\le T\}}g(\tau)]$ (paths are
continuous at a fixed time, so $V_{t-}=V_t$ a.s.), and $\tau$ has density $h(t)S(t)$. For the approximation, put
$\mathrm{EE}^*(t)=e^{-rt}\mathrm{EPE}$ and $hS(t)=he^{-ht}$, then use $(1-R)h=s$.
\end{proof}
The wider the spread (default more likely) and the longer the duration (more time to default), the larger the
charge \ts{13}{26}{14:55}.

\begin{example}[Size of CVA]
A 5-year uncollateralised swap book with $\mathrm{EPE}=\$5$M, counterparty spread $s=100$\,bp, $R=40\%$ (so
$h=1.67\%$) and $r=3\%$: $D=(1-e^{-0.0467\cdot5})/0.0467=4.46$ and $\mathrm{CVA}\approx-5\cdot0.01\cdot4.46
=-\$0.22$M. If the spread doubles, the CVA roughly doubles: a mark-to-market loss with no default at all.
\end{example}

\begin{intuition}[CVA is an option on a basket]
$(V_\tau-C_\tau)^+$ is the payoff of a call on the netting set with strike equal to the collateral, exercised at
the random time $\tau$. CVA is a default-probability-weighted strip of such options. Since a counterparty trades
many products (rates, FX, credit, equity, commodities, mortgages), it is an option on a \emph{basket of cross-asset
derivatives} \ts{13}{26}{20:11}. Pricing it needs a joint simulation of all risk factors of all trades: this is
what Tang calls enterprise-level modelling (\cref{ch08:sec-martingale}), and the reason his group at Morgan Stanley
was called Cross Asset Modeling.
\end{intuition}

\begin{finance}[CVA in the bank]
CVA changes the true economic price of a trade, so it is charged to clients and hedged dynamically by a CVA
desk, which buys CDS protection and hedges the market factors that drive $V$. It also enters regulatory capital:
Basel~III added a CVA capital charge, so trades with large CVA consume more capital and lower the return on
equity \ts{13}{26}{5:17}. In 2008, CVA mark-to-market losses were larger than the losses from actual defaults
\ts{13}{26}{6:22}.
\end{finance}

\begin{exercise}[CVA from exposure (not from the course)]
\leavevmode
\begin{enumerate}[(a)]
  \item Derive \cref{ch08:prop-cvaee} and the approximation \eqref{ch08:eq-cvabote}.
  \item For the exposure profile
      $\mathrm{EE}(t)=\sigma\sqrt t(1-t/T)/\sqrt{2\pi}$ of \cref{ch08:fig-exposure}, with $r=0$ and constant hazard rate $h$,
      compute CVA to first order in $h$ and compare with $\mathrm{EPE}\cdot s\cdot T$.
  \item Explain qualitatively how wrong-way
      risk changes the result.
\end{enumerate}
\end{exercise}

\subsection{DVA and first-to-default}
The same logic applied to \emph{our own} default gives a \emph{liability} CVA \ts{13}{26}{21:23}. With $\bar\tau$
our default time and $\bar R$ our recovery, the slides write
\[
  \mathrm{CVA}_{\mathrm{MNE}} = -\E_0^{\Q}\!\left[\ind_{\{\bar\tau\le T\}}\,
      \frac{(V_{\bar\tau-}-C_{\bar\tau-})^-\,(1-\bar R_{\bar\tau})}{\beta_{\bar\tau}}\right]
  = \E_0^{\Q}\!\left[\ind_{\{\bar\tau\le T\}}\,
      \frac{(C_{\bar\tau-}-V_{\bar\tau-})^+\,(1-\bar R_{\bar\tau})}{\beta_{\bar\tau}}\right]\ \ge 0 .
\]
This is a \emph{benefit}: if we default while owing money, we pay only a fraction of it. It is usually called
\emph{debit valuation adjustment} (DVA) \ts{13}{26}{22:27}.

\begin{coursenote}[Sign convention on slide~6 of \file{lecnote26}]
For the displayed formula to be a benefit, $(x)^-$ must be read as $\min(x,0)\le0$, as Tang explains
(``the positive sign becomes a negative sign''). With the other common convention $x^-=\max(-x,0)$ the
formula would carry the opposite sign. The right-hand side above avoids the ambiguity.
\end{coursenote}

Should the two adjustments know about each other? If we default first, the counterparty's later default no longer
matters to us, and vice versa. A \emph{bilateral} CVA puts first-to-default indicators $\ind_{\{\tau\le\min(\bar\tau,T)\}}$
and $\ind_{\{\bar\tau\le\min(\tau,T)\}}$ into the two terms (not written in class). Tang's view is that for pricing
this refinement does not need to be modelled; he gave references on the slides rather than discussing it
\ts{13}{26}{23:30}.

\begin{finance}[Why DVA is controversial (not discussed in class)]
DVA means that a bank books a gain when its own credit worsens, and it can hardly be hedged, since a bank cannot
sell protection on itself. Regulators therefore do not let DVA gains count towards capital, even though
accounting standards include DVA in fair value.
\end{finance}

\subsection{Funding and liquidity risk}
Second example \ts{13}{26}{23:30}: same trade, value $+\$100$M, uncollateralised, with market risk and counterparty
risk properly hedged. Is anything left? Students proposed interest-rate risk (hedged by assumption) and
``key-man'' risk \ts{13}{26}{25:40}. Tang's answer is \emph{funding liquidity risk} \ts{13}{26}{25:40}:
\begin{itemize}
  \item hedges are typically futures or collateralised trades. When our uncollateralised receivable is worth
        $+100$ the hedge is worth $-100$, and we must post $100$ of cash \emph{now}, while the receivable
        arrives only later. That cash must be borrowed at our unsecured funding spread \ts{13}{26}{26:47};
  \item the amount depends on market moves (\emph{contingent} liquidity risk), and failing to raise it is how
        firms die: ``you may end up like Lehman'' \ts{13}{26}{27:51};
  \item symmetrically, an uncollateralised \emph{payable} generates a funding \emph{benefit} (we receive
        collateral on its hedge), which can partially offset the funding need of the receivables
        \ts{13}{26}{27:51};
  \item further risks: unexpected tail events and (equity) capital \ts{13}{26}{28:54}.
\end{itemize}
In the spirit of \eqref{ch08:eq-cvabote}, a first estimate of the funding cost is $-s_F\times\mathrm{EPE}\times D$
and of the benefit $+s_F\times|\mathrm{ENE}|\times D$, with $s_F$ our own funding spread (a common simplification,
not from the lecture).

\begin{finance}[The XVA family (terminology not used in class)]
The adjustments to the default-free price are nowadays collectively called XVA: CVA (counterparty default),
DVA (own default), FVA (funding of uncollateralised positions, split into cost FCA and benefit FBA), KVA (cost
of regulatory capital held against the trade, Tang's ``capital''), and MVA (cost of funding the \emph{initial
margin} of \cref{ch08:sec-margin} over the life of the trade). The 2013 lecture covered CVA, DVA, funding and
capital. The 2024 lecture is about the world that the margin rules created, in which CVA is smaller because
trades are collateralised and the cost has moved into margin; its reading list cites J.~Gregory, \emph{The xVA
Challenge}.
\end{finance}

\subsection{Wrong-way risk}\label{ch08:sec-wwr}
\begin{definition}[Wrong-way risk]
\emph{Wrong-way risk} (WWR) is present when exposure to a counterparty is positively correlated with its
probability of default: the counterparty owes us more exactly when it is more likely to default
\ts{13}{26}{33:02}. (\emph{Right-way} risk is the opposite.) Then the independence assumption of
\cref{ch08:prop-cvaee} fails, and the relevant exposure is the \emph{conditional} one,
$\mathrm{CVA}=-(1-R)\int_0^T\E^{\Q}\bigl[(V_t-C_t)^+/\beta_t\,\big|\,\tau=t\bigr]\dd\Q(\tau\le t)$, which is
larger than $\mathrm{EE}^*(t)$ under WWR.
\end{definition}

\begin{casestudy}[Berkshire Hathaway's index puts (2013 slides 8--9)]
\textbf{The trade} \ts{13}{26}{29:54}. Selling a put below the current price is an alternative to buying the
stock: similar downside, no upside beyond the premium, and a steady income if repeated (``name your price and
get paid for it''). Around 2005--2006, as Tang recalls, Berkshire Hathaway sold long-dated (10--15 year) puts on
four indices (S\&P~500, FTSE~100, Euro Stoxx~50, Nikkei~225) and collected about $\$4$bn of premium \emph{without
posting collateral}. It was one of the largest derivative cash outflows of its time for the dealers, who were the
buyers \ts{13}{26}{30:55}. Had Berkshire posted collateral, the 2008 sell-off would have required many billions of
it (see page~16 of Berkshire's 2012 shareholder letter, cited on the slides).

\textbf{The dealer's risks} \ts{13}{26}{31:56}. Handling long-dated equity options was not the difficulty; many dealers
could not handle the \emph{counterparty} side: CVA on a large uncollateralised receivable; \emph{wrong-way} CVA,
because in a sell-off the puts gain value while an insurer with large equity holdings becomes more likely to
default \ts{13}{26}{33:02}; and \emph{wrong-way funding}, because the dealer paid $\$4$bn up front and had to fund
its collateralised hedges exactly in a crisis \ts{13}{26}{34:06}. The dealer therefore charged CVA, wrong-way CVA,
funding and wrong-way funding costs. It then had to risk-manage them: a static CDS hedge (fine for a bond) leaves
\emph{gap risk} because the exposure changes over time \ts{13}{26}{35:11}. Part of the risk was sold to investors
in Berkshire's bonds through a \emph{credit-linked note} paying a higher coupon, and through tranched portfolio
protection \ts{13}{26}{36:17}.

\textbf{Berkshire's side} \ts{13}{26}{37:18}. Selling puts is selling insurance on the equity market; the premium
is day-one cash, a cheap source of funding (Buffett compared it to borrowing at about 1\%); and there is no cash
outflow until maturity. So Berkshire suffered mark-to-market losses in 2008 but no liquidity drain.

\textbf{Feedback} \ts{13}{26}{38:25}. In the crisis, dealers' CVA hedging meant buying more and more CDS
protection on Berkshire, which widened Berkshire's CDS spread beyond what fundamentals suggested. Bond investors
then demanded higher coupons, raising Berkshire's funding cost, so Berkshire approached dealers about posting
collateral after all: ``no free lunch'' \ts{13}{26}{39:30}.

\textbf{Who loses?} Tang's answer: the end buyers of the puts, the hedgers who pay the premium and gain nothing if
the market does not fall far enough \ts{13}{26}{40:33}.
\end{casestudy}

\subsection{The CVA conundrum}
To hedge CVA on counterparty A you buy credit protection on A from counterparty B. Now you have counterparty risk
to B, so you buy protection on B from C, and so on: an infinite series \ts{13}{26}{41:41}. This also undermines
the theory: arbitrage pricing is replication with hedge instruments, and here replication needs infinitely many
of them \ts{13}{26}{43:52}. The practical way to stop the series at the first term is to buy the protection on A
\emph{fully collateralised}, typically from a dealer under a CSA: collateral is settled as the value moves, much
like daily settlement of futures, so the residual risk to B is negligible \ts{13}{26}{43:52}.

\begin{coursenote}[Slides 21--25 of \file{lecnote26}: the extinguisher (not presented)]\label{ch08:rem-extinguisher}
Tang ran out of time and referred students to the slides \ts{13}{26}{1:16:34}. The question: normally a dealer
\emph{charges} CVA; how can it structure a trade in which it \emph{pays} the counterparty for the latter's default
risk? Answer: a trade with positive value to the dealer that increases with the counterparty's default probability,
namely an \emph{extinguisher} (zero-recovery) swap, which knocks out at zero value if either party defaults. It is
worth something when the dealer expects to have more future \emph{payables} than receivables. Example: an
emerging-market sovereign issues a USD bond and swaps it into local currency with a dealer (receives USD, pays
local currency). Because of the interest-rate differential (the FX forward), the correlation between FX and
credit, and the jump of the FX rate at default, the dealer expects to owe money on the swap; the extinguisher
cancels this debt at the sovereign's default. In effect the dealer has bought protection on the sovereign from the
sovereign itself, and can sell it on to a third party. Issues: it reduces recovery for the other creditors; the
default event must reference the underlying bond; there are legal and franchise risks; and why would the sovereign
default, knowing it would lose its receivable on the swap? Slide~25 places all this in the ``bigger picture'' of a
bank's objectives: revenue, risk management (market, counterparty, liquidity, capital) and return on capital,
with the non-linear portfolio effects requiring enterprise-level models.
\end{coursenote}

% ==================================================================
\section{Risk measures: value at risk and expected shortfall}\label{ch08:sec-riskmeasures}
Initial margin is a risk measure of a portfolio's value change, and to optimise it we must know which risk
measures are well behaved. This section follows the first half of the 2024 lecture.

\subsection{Coherent risk measures}
A risk measure should capture both the \emph{size} of potential losses and their \emph{probability}
\ts{24}{10}{8:25}. Below $L$ denotes the \emph{loss} of a portfolio over a horizon (a random variable; gains are
negative losses) and $\rho(L)\in\R$ its risk.

\begin{definition}[Coherent risk measure, Artzner et al.\ 1999]\label{ch08:def-coherent}
$\rho$ is \emph{coherent} if for all losses $L,L_1,L_2$:
\begin{enumerate}[(i)]
  \item normalised: $\rho(0)=0$;
  \item positively homogeneous: $\rho(\lambda L)=\lambda\rho(L)$ for $\lambda\ge0$;
  \item translation invariant: $\rho(L+a)=\rho(L)+a$ for a deterministic amount $a\in\R$;
  \item monotone: $L_1\le L_2$ in all scenarios $\Rightarrow\rho(L_1)\le\rho(L_2)$;
  \item subadditive: $\rho(L_1+L_2)\le\rho(L_1)+\rho(L_2)$ \ts{24}{10}{9:29}.
\end{enumerate}
\end{definition}
Subadditivity is the important one: merging two portfolios should not create risk, and it may reduce it through
diversification \ts{24}{10}{10:31}.

\begin{coursenote}[Erratum: slide~6 of \file{lec10}]
The slide states homogeneity for all $\alpha$; it holds only for $\alpha\ge0$ (doubling a position doubles its
risk, but reversing it does not negate the risk). It also writes translation invariance as
$\rho(X+A)=\rho(X)+\alpha$ for ``a deterministic portfolio with guaranteed \emph{return} $\alpha$''. With $X$ a loss,
as in the rest of the slides, adding a guaranteed \emph{gain} $\alpha$ \emph{lowers} the risk:
$\rho(X-\alpha)=\rho(X)-\alpha$; equivalently, adding a deterministic \emph{loss} $a$ raises it by $a$, as in
(iii) above. Finally, the reference list dates Artzner--Delbaen--Eber--Heath to 2022; the paper appeared in
\emph{Mathematical Finance} in 1999, the year quoted in class.
\end{coursenote}

\subsection{VaR and ES}
\begin{definition}[VaR and ES]\label{ch08:def-vares}
For a confidence level $\beta\in(0,1)$ and a horizon (implicit in $L$):
\begin{align*}
  \mathrm{VaR}_\beta(L) &= \inf\{\ell\in\R:\ \Prob(L\le\ell)\ge\beta\}, \\
  \mathrm{ES}_\beta(L)  &= \frac{1}{1-\beta}\int_\beta^1\mathrm{VaR}_u(L)\dd u .
\end{align*}
If $L$ has a continuous density $f_L$ then $\int_{-\infty}^{\mathrm{VaR}_\beta}f_L=\beta$ and
\[
  \mathrm{ES}_\beta(L)=\E\bigl[L\,\big|\,L\ge\mathrm{VaR}_\beta\bigr]
  =\frac{1}{1-\beta}\int_{\mathrm{VaR}_\beta}^\infty x\,f_L(x)\dd x .
\]
\end{definition}
VaR answers ``what is the loss that I am $\beta$ sure not to exceed over the next 10 days?'' (e.g.\ $\$10$M at
$\beta=99\%$); ES answers ``if I am in the worst $1-\beta$ of cases, how much do I lose on average?''
\ts{24}{10}{10:31}. Clearly $\mathrm{ES}_\beta\ge\mathrm{VaR}_\beta$, since ES averages quantiles $\mathrm{VaR}_u\ge
\mathrm{VaR}_\beta$ \ts{24}{10}{12:35}. VaR was popularised by JPMorgan in the 1990s and ES appeared around 2000;
both are used heavily by regulators and in counterparty risk.

\begin{figure}[ht]
\centering
\begin{tikzpicture}
\begin{axis}[width=0.85\linewidth,height=5.2cm,xmin=-4,xmax=6.5,ymin=0,ymax=0.43,
  xlabel={loss $L$},ylabel={density},ytick=\empty,
  tick label style={font=\scriptsize},label style={font=\small},
  declare function={fb(\x)=0.97*exp(-\x*\x/2)/sqrt(2*pi)+0.03*exp(-(\x-5)*(\x-5)/0.18)/(0.3*sqrt(2*pi));
                    fr(\x)=exp(-\x*\x/(2*1.53983))/(1.24092*sqrt(2*pi));}]
\addplot[fill=ideacol,fill opacity=0.25,draw=none,domain=2.041:6.5,samples=120]{fb(x)}\closedcycle;
\addplot[fill=thmcol,fill opacity=0.25,draw=none,domain=2.041:6.5,samples=80]{fr(x)}\closedcycle;
\addplot[ideacol,thick,domain=-4:6.5,samples=200]{fb(x)};
\addplot[thmcol,thick,domain=-4:6.5,samples=200]{fr(x)};
\draw[dashed] (axis cs:2.041,0) -- (axis cs:2.041,0.30) node[above,font=\scriptsize]{$\mathrm{VaR}_{95\%}=2.04$};
\draw[thmcol,thick,->] (axis cs:2.56,0.13) node[above,font=\scriptsize]{ES $2.56$} -- (axis cs:2.56,0.005);
\draw[ideacol,thick,->] (axis cs:3.96,0.13) node[above,font=\scriptsize]{ES $3.96$} -- (axis cs:3.96,0.005);
\end{axis}
\end{tikzpicture}
\caption{Same VaR, different tails (in the spirit of slide~8 of \file{lec10}; numbers not from the course). Red:
$N(0,1.24^2)$. Blue: $97\%\ N(0,1)+3\%\ N(5,0.3^2)$, a book that usually makes small gains but occasionally
loses a lot. Both have $\mathrm{VaR}_{95\%}=2.04$, but the blue $\mathrm{ES}_{95\%}$ is $3.96$ against $2.56$:
VaR does not see what lies beyond the quantile.}\label{ch08:fig-fattail}
\end{figure}

\paragraph{Fat tails.} VaR is blind to the shape of the tail beyond the quantile (\cref{ch08:fig-fattail})
\ts{24}{10}{12:35}. Such a distribution is not artificial: a desk that sells out-of-the-money options collects a
small premium most of the time and occasionally suffers a large loss when the options end in the money
\ts{24}{10}{14:37}.

\paragraph{VaR is not subadditive.} Slide~9 \ts{24}{10}{14:37}: two independent portfolios A and B each lose $100$
with probability $4\%$ and $0$ otherwise. For each, $\mathrm{VaR}_{0.95}=0$ (the $4\%$ tail lies beyond the
quantile) and $\mathrm{ES}_{0.95}=(4\cdot100+1\cdot0)/5=80$. The sum A+B loses $200$ w.p.\ $0.16\%$, $100$ w.p.\
$7.68\%$, $0$ w.p.\ $92.16\%$, so $\mathrm{VaR}_{0.95}=100>0+0$ while $\mathrm{ES}_{0.95}=(0.16\cdot200+4.84\cdot100)/5
=103.2\le160$ \ts{24}{10}{15:39}. ES is subadditive in general (it is coherent; the slides point to a paper with
seven different proofs), but proving it takes some work. Conversely, for normal-like (elliptical) distributions
VaR \emph{is} subadditive, so in practice the failure is less dramatic than such contrived examples suggest
\ts{24}{10}{16:45}. For optimisation, however, it is fatal, because of the following fact.

\begin{proposition}[Subadditivity and convexity]\label{ch08:prop-convex}
Let $\rho$ be positively homogeneous. Then $\rho$ is subadditive iff it is convex.
\end{proposition}
\begin{proof}
If $\rho$ is subadditive, then for $\lambda\in[0,1]$:
$\rho(\lambda L_1+(1-\lambda)L_2)\le\rho(\lambda L_1)+\rho((1-\lambda)L_2)=\lambda\rho(L_1)+(1-\lambda)\rho(L_2)$.
If $\rho$ is convex: $\rho(L_1+L_2)=2\rho\bigl(\tfrac12L_1+\tfrac12L_2\bigr)\le\rho(L_1)+\rho(L_2)$.
\end{proof}
Without homogeneity the equivalence fails: $\sqrt{|x|}$ is subadditive but not convex. The slides (slide~33) state
``convexity is equivalent to sub-additivity''; the proposition makes precise when this is true.
Since VaR is not subadditive, it is not convex as a function of the portfolio, and non-convex functions are much
harder to minimise than convex ones \ts{24}{10}{16:45}.

\paragraph{Spectral risk measures.} A risk measure can be described by the weight it puts on each quantile,
$M_\phi(L)=\int_0^1\phi(u)\,\mathrm{VaR}_u(L)\dd u$ with $\phi\ge0$, $\int\phi=1$ \ts{24}{10}{17:45}. VaR puts all
its weight on $u=\beta$ (a delta function); ES puts a flat weight $1/(1-\beta)$ on $u>\beta$; an
\emph{exponential} spectral measure uses $\phi(u)\propto e^{-(1-u)/\gamma}$, weighting extreme losses most.
$M_\phi$ is coherent iff $\phi$ is non-decreasing (Acerbi, 2002): a sensible risk-averse measure never gives
worse losses less weight. This fails for VaR, whose weight drops to zero beyond $\beta$.

\begin{intuition}[Quantile versus conditional mean]
VaR is a quantile, like a median: it does not move if you drag the extreme tail further out. ES is the first moment
of the tail, so it does. Convexity is the property that makes a minimisation problem a bowl with one bottom (the
same role convexity of the free energy plays for thermodynamic stability); a risk measure without it can have
many local minima in portfolio space.
\end{intuition}

\begin{exercise}[VaR is not subadditive (not from the course)]
\leavevmode
\begin{enumerate}[(a)]
  \item Verify all the numbers of the A, B, A+B example (VaR and ES at $95\%$).
  \item Replace the $4\%$ loss probability by
      $p$ and find the range of $p$ for which $\mathrm{VaR}_{0.95}(A+B)>\mathrm{VaR}_{0.95}(A)+\mathrm{VaR}_{0.95}(B)$.
  \item Show that for jointly normal losses VaR is subadditive at any $\beta\ge\frac12$.
\end{enumerate}
\end{exercise}

\subsection{Normal model and time scaling}
\begin{proposition}[Normal VaR and ES]\label{ch08:prop-normal}
If $L\sim N(\mu,\sigma^2)$ and $z_\beta=N^{-1}(\beta)$, with $\varphi$ the standard normal density, then
\[
  \mathrm{VaR}_\beta=\mu+\sigma z_\beta,\qquad \mathrm{ES}_\beta=\mu+\sigma\,\frac{\varphi(z_\beta)}{1-\beta}
  =\mu+\sigma\,\frac{e^{-z_\beta^2/2}}{\sqrt{2\pi}\,(1-\beta)}\ts{24}{10}{18:45}.
\]
\end{proposition}
\begin{proof}
$L=\mu+\sigma Z$ and quantiles are equivariant under increasing affine maps. For ES,
$\E[Z\mid Z\ge z]=\frac{1}{1-\Phi(z)}\int_z^\infty x\varphi(x)\dd x=\frac{\varphi(z)}{1-\Phi(z)}$, since
$x\varphi(x)=-\varphi'(x)$.
\end{proof}
Numerically, $z_{0.99}=2.326$ and $\varphi(z_{0.99})/0.01=2.665$, while $\varphi(z_{0.975})/0.025=2.338$. Hence
under normality the $99\%$ VaR ($\mu+2.326\sigma$) and the $97.5\%$ ES ($\mu+2.338\sigma$) are almost the same
number: this is how the regulators' switch from 99\% VaR to 97.5\% ES was calibrated (slide~17).

\paragraph{Square root of time.} If daily changes $\Delta P_i$ are i.i.d.\ $N(0,\sigma^2)$, the $T$-day change has
standard deviation $\sigma\sqrt T$, so $T$-day VaR $=\sqrt T\times$ one-day VaR, and the same for ES
\ts{24}{10}{19:47}. Real returns are autocorrelated: a bad day tends to be followed by another.

\begin{proposition}[Autocorrelated daily changes]\label{ch08:prop-autocorr}
If $(\Delta P_i)$ is stationary with variance $\sigma^2$ and $\Corr(\Delta P_i,\Delta P_{i+k})=\rho^k$ (an AR(1)
structure), then
\[
  \Var\Bigl(\sum_{i=1}^T\Delta P_i\Bigr)=\sigma^2\Bigl[T+2\sum_{k=1}^{T-1}(T-k)\rho^k\Bigr],
\]
and for normal changes with zero mean the ratio of $T$-day to one-day VaR (or ES) is the square root of the bracket.
For $T=2$ it is $\sqrt{2(1+\rho)}$, as on slide~11.
\end{proposition}
\begin{proof}
Expand the variance of the sum: there are $T$ diagonal terms and $2(T-k)$ pairs at lag $k$.
\end{proof}
\begin{center}\small
\begin{tabular}{@{}lcccccc@{}}
\toprule
 & $T=1$ & $2$ & $5$ & $10$ & $50$ & $250$\\
\midrule
$\rho=0$    & 1.00 & 1.41 & 2.24 & 3.16 & 7.07 & 15.81\\
$\rho=0.05$ & 1.00 & 1.45 & 2.33 & 3.31 & 7.43 & 16.62\\
$\rho=0.1$  & 1.00 & 1.48 & 2.42 & 3.46 & 7.80 & 17.47\\
$\rho=0.2$  & 1.00 & 1.55 & 2.62 & 3.79 & 8.62 & 19.35\\
\bottomrule
\end{tabular}
\end{center}
The table of slide~11 is reproduced exactly by the formula above. For the EURO STOXX the observed lag-one
autocorrelation is about $0.1$ (between $0.05$ and $0.15$ for other indices) \ts{24}{10}{20:54}, so the usual
10-day rule $\sqrt{10}=3.16$ underestimates the true factor $3.46$ by about $10\%$; yet almost everybody uses
$\sqrt T$ \ts{24}{10}{21:59}. (The observed higher-lag autocorrelations on the slide are close to $0$ rather than
$\rho^k$; with only a lag-one correlation the factor is $\sqrt{10+2\cdot9\cdot0.1}=3.44$, practically the same.)

\begin{exercise}[Normal model and FRTB (not from the course)]
\leavevmode
\begin{enumerate}[(a)]
  \item Prove \cref{ch08:prop-normal}.
  \item Find the confidence level $\beta'$ at which the normal ES equals the normal
      $99\%$ VaR, and compare it with the $97.5\%$ used by FRTB.
  \item With $\bm\delta=(3000,4000,2000,1000)$ and the covariance
      matrix of slide~16 (rows $(1.02,0.97,0.80,0.94)$, $(0.97,1.08,0.82,0.96)$, $(0.80,0.82,0.89,0.76)$,
      $(0.94,0.96,0.76,0.96)$, all $\times10^{-4}$), reproduce $\sigma_P=96.6$, $\mathrm{VaR}=225$, $\mathrm{ES}=257$, and
      compute the 10-day values with and without the autocorrelation correction for $\rho=0.1$.
\end{enumerate}
\end{exercise}

\subsection{Estimating VaR and ES}\label{ch08:sec-estimation}

\begin{casestudy}[Historical simulation on four European indices (slides 12--15)]
\textbf{Data} \ts{24}{10}{21:59}. Daily closes of IBEX~35 (Spain), CAC~40 (France), DAX (Germany) and EURO STOXX~50,
2006--2023, containing two stress periods (the 2008 crisis and COVID in 2020). All four are in euros, so there is no FX
conversion; holidays differ across countries: days when all markets were closed were dropped, otherwise the
missing prices were forward-filled, which creates zero returns \ts{24}{10}{23:01}. (Shepherd's aside: in the
industry a surprising share of time goes into holiday calendars and FX.)

\textbf{Method} \ts{24}{10}{24:02}. With $v_i^j$ the level of index $j$ on day $i$ and positions
$\bm\delta=(3000,4000,2000,1000)$ in CAC, DAX, IBEX, EURO STOXX, the P\&L of day $i$ is
$\Delta P_i=\sum_j\delta_j\,(v_i^j-v_{i-1}^j)/v_{i-1}^j$. For a window of $n=500$ days (two years), sort the
losses: the $99\%$ VaR is the 5th worst ($1\%$ of $500$) and the $99\%$ ES the mean of the 4 worst; the $97.5\%$ VaR
is the 12th worst and the ES the mean of the 11 worst \ts{24}{10}{27:07}. (Some include the VaR scenario itself in
the average; the market convention is to exclude it \ts{24}{10}{28:10}.)

\textbf{Results} (slides 14--15; checked against the listed 12 largest losses):
\begin{center}\small
\begin{tabular}{@{}lcccc@{}}
\toprule
window & $\mathrm{VaR}_{99}$ & $\mathrm{ES}_{99}$ & $\mathrm{VaR}_{97.5}$ & $\mathrm{ES}_{97.5}$\\
\midrule
2022--2023 (recent)   & 320 & 377 & 227 & 303\\
2020--2021 (COVID)    & 450 & 815 & 306 & 534\\
2015--2016 (calm)     & 297 & 477 & 246 & 348\\
2007--2008 (crisis)   & 448 & 689 & 370 & 505\\
\bottomrule
\end{tabular}
\end{center}
\textbf{Lesson.} The choice of the historical window matters far more than any other modelling choice
\ts{24}{10}{26:07}. This is why regulators added \emph{stressed VaR}: besides the most recent window, VaR is also
computed on a fixed period known to be stressful (2008 or COVID) \ts{24}{10}{27:07}.
\end{casestudy}

\paragraph{Model building (variance--covariance).} Assume index returns $\Delta x_i\sim N(0,\sigma_i^2)$ with
correlations $\rho_{ij}$, covariance $C_{ij}=\rho_{ij}\sigma_i\sigma_j$, and a linear portfolio
$\Delta P=\sum_i\delta_i\Delta x_i$ with sensitivities $\delta_i=\partial P/\partial x_i$ \ts{24}{10}{28:10}. Then
$\Delta P\sim N(0,\sigma_P^2)$ with
\[
  \sigma_P^2=\sum_{i,j}\rho_{ij}\delta_i\delta_j\sigma_i\sigma_j=\bm\delta\T C\bm\delta,\qquad
  \mathrm{VaR}(T,\beta)=\sqrt T\,z_\beta\,\sigma_P,\qquad
  \mathrm{ES}(T,\beta)=\sqrt T\,\frac{\varphi(z_\beta)}{1-\beta}\,\sigma_P .
\]
With the 2022--2023 daily covariance matrix of slide~16 (entries about $10^{-4}$, i.e.\ daily volatilities of
about $1\%$) one gets $\sigma_P=96.6$, $\mathrm{VaR}(1,99\%)=225$ and $\mathrm{ES}(1,99\%)=257$ \ts{24}{10}{29:13},
well below the historical $320$ and $377$ for the same period: the empirical distribution has fatter tails than the
normal \ts{24}{10}{30:15}.

\begin{coursenote}[Erratum: ES formula on slide~16 of \file{lec10}]
The slide writes $\mathrm{ES}(T,\beta)=\sqrt T\,e^{-N^{-1}(\beta)}\sigma_P/(\sqrt{2\pi}(1-\beta))$. The exponent
must be $-[N^{-1}(\beta)]^2/2$, as on slide~11 and in \cref{ch08:prop-normal}. The quoted $\mathrm{ES}=257$ was
computed with the correct formula ($2.665\times96.6$); the misprinted one would give $376$.
\end{coursenote}

\subsection{Back-testing}
A VaR model can be tested by counting \emph{exceptions}, days on which the realised loss exceeded the VaR
\ts{24}{10}{30:15}. If the model is right, exceptions are independent Bernoulli trials with probability
$p=1-\beta$, so the number $N$ of exceptions in $n$ days is $\mathrm{Bin}(n,p)$ and
\[
  \Prob(N\ge m)=\sum_{k=m}^n\binom nk p^k(1-p)^{n-k}.
\]
The rule of thumb is to reject the model if this probability is below $5\%$ \ts{24}{10}{31:18}. In the example the
normal $\mathrm{VaR}(1,99\%)=225$ was exceeded on $12$ of the $500$ days, while $5$ were expected.

\begin{coursenote}[Erratum: the back-test on slide~16 and in class]
The slide's formula has $(1-p)^k$; the binomial probability needs $(1-p)^{n-k}$. More importantly, for
$n=500$, $m=12$, $p=0.01$ the correct value is $\Prob(N\ge12)=0.52\%$, not $5.2\%$. By the stated $5\%$ rule the
normal model is therefore \emph{rejected} (even at the $1\%$ level), whereas in class it was accepted as ``not
overly aggressive'' \ts{24}{10}{1:20:34}. At the $5\%$ level the model is accepted with up to $9$ exceptions
($\Prob(N\ge9)=6.7\%$) and rejected from $10$ on ($\Prob(N\ge10)=3.1\%$). The rejection also agrees with the
comparison with the historical VaR ($225$ against $320$). Note that this back-test is in-sample (covariance and
exceptions from the same window), which if anything favours the model.
\end{coursenote}

Back-testing VaR is easy; back-testing ES is not, and this is the main reason VaR is still used despite its flaws
\ts{24}{10}{32:23}. ES is a property of the whole tail of a distribution, while each day provides a single draw: one
realised loss cannot be compared with a tail mean without assumptions on the shape of the tail
\ts{24}{10}{1:17:27}. Regulation moved from VaR (Basel~I, 10-day $99\%$) to stressed VaR (Basel~2.5) to ES at
$97.5\%$ with variable liquidity horizons (FRTB, the Fundamental Review of the Trading Book). Banks still
back-test the VaR: they build one distributional model, compute capital from its ES, and accept the ES if the VaR of
the same model back-tests well \ts{24}{10}{1:16:23}. The ``model'' is the distributional assumption; VaR and ES are
two numbers read off it \ts{24}{10}{1:18:28}.

\begin{exercise}[Back-testing (not from the course)]
\leavevmode
\begin{enumerate}[(a)]
  \item Compute $\Prob(N\ge12)$ for $N\sim\mathrm{Bin}(500,0.01)$ exactly and with the Poisson approximation.
  \item Find the
      largest number of exceptions accepted at the $5\%$ level.
  \item Why is the in-sample back-test of the example biased in
      favour of the model?
  \item Explain why a single observed loss cannot ``test'' an ES forecast, and propose a test that
      uses the ES forecasts and realised losses of many days.
\end{enumerate}
\end{exercise}

\subsection{Expected shortfall as an optimisation problem}
Let $x\in X\subset\R^n$ be a portfolio, $y\in\R^m$ the random market change with density $p(y)$, and $f(x,y)$ the
loss. Write $\psi(x,\alpha)=\Prob(f(x,y)\le\alpha)$, so $\mathrm{VaR}_\beta(x)=\min\{\alpha:\psi(x,\alpha)\ge\beta\}$.

\begin{theorem}[Rockafellar--Uryasev]\label{ch08:thm-ru}
Let
\[
  F_\beta(x,\alpha)=\alpha+\frac{1}{1-\beta}\int_{\R^m}\bigl[f(x,y)-\alpha\bigr]^+p(y)\dd y .
\]
If $\psi(x,\cdot)$ is continuous, then $F_\beta(x,\cdot)$ is convex and continuously differentiable,
$\mathrm{VaR}_\beta(x)$ is a minimiser, and $\min_\alpha F_\beta(x,\alpha)=\mathrm{ES}_\beta(x)$. Consequently
\[
  \min_{x\in X}\mathrm{ES}_\beta(x)=\min_{(x,\alpha)\in X\times\R}F_\beta(x,\alpha),
\]
and if $f(\cdot,y)$ is convex and $X$ convex, this is a convex problem \ts{24}{10}{33:27}.
\end{theorem}
\begin{proof}
$\alpha\mapsto[f-\alpha]^+$ is convex, so $F_\beta$ is convex in $\alpha$, and
\[
  \frac{\partial F_\beta}{\partial\alpha}=1-\frac{\Prob(f>\alpha)}{1-\beta}=1+\frac{\psi(x,\alpha)-1}{1-\beta}
  =\frac{\psi(x,\alpha)-\beta}{1-\beta},
\]
which vanishes at $\alpha_\beta=\mathrm{VaR}_\beta(x)$. There
\[
  \int[f-\alpha_\beta]^+p\dd y=\int_{f\ge\alpha_\beta}f\,p\dd y-\alpha_\beta\bigl(1-\psi(x,\alpha_\beta)\bigr)
  =(1-\beta)\mathrm{ES}_\beta-\alpha_\beta(1-\beta),
\]
so $F_\beta(x,\alpha_\beta)=\alpha_\beta+\mathrm{ES}_\beta-\alpha_\beta=\mathrm{ES}_\beta$ (slides 18--20). If $f$ is jointly
convex in $(x,\alpha)$, so is $[f-\alpha]^+$ and hence $F_\beta$; minimising first over $\alpha$ and then over $x$ is
the same as minimising jointly.
\end{proof}
With $S$ equally likely scenarios $y_s$ (historical or simulated) and linear losses $f(x,y_s)=-x\T r_s$, the problem
becomes a \emph{linear program}:
\[
  \min_{x,\alpha,u}\ \alpha+\frac{1}{(1-\beta)S}\sum_{s=1}^Su_s\quad\text{s.t.}\quad u_s\ge -x\T r_s-\alpha,\ \ u_s\ge0,
  \ \ x\in X ,
\]
``basically a linear problem'' \ts{24}{10}{34:27}. No such representation exists for VaR. So ES is harder to
back-test but much easier to optimise, and for margin optimisation the latter matters more.

\begin{exercise}[ES as a linear program (not from the course)]
\leavevmode
\begin{enumerate}[(a)]
  \item Complete the proof of \cref{ch08:thm-ru} for a discrete distribution with $S$ equally likely scenarios, taking
      $\beta S$ to be an integer, and check that $\min_\alpha F$ is the average of the $(1-\beta)S$ largest losses.
  \item Write the LP that minimises the $97.5\%$ ES of a long-only fully invested portfolio of the four indices over the
      500 historical scenarios of 2022--2023.
  \item Why can the same trick not be used for VaR?
\end{enumerate}
\end{exercise}

% ==================================================================
\section{Margin: variation margin, initial margin and ISDA SIMM}\label{ch08:sec-margin}

\subsection{Why margin}
Consider a bank A that is long a swap with C and short an offsetting swap with B (slide~21). A receives cash flows
$\{T_i\}$ and pays $\{S_i\}$ on each swap, typically every three or six months; before the next payment date the
counterparty may default and not pay \ts{24}{10}{35:33}. To mitigate this the parties exchange \emph{variation
margin} daily, equal to the current value $V$ of the trades \ts{24}{10}{36:34}.

If C defaults it stops paying VM. A must replace or unwind the trade with C, which takes some days; meanwhile its
hedge with B keeps moving and A keeps paying VM to B. If the value of the defaulted trade to A rises in the
meantime, A has to pay more for the replacement than the collateral it holds, and that is a loss
\ts{24}{10}{37:35}. \emph{Initial margin} covers this loss.

\begin{definition}[Margin]
\emph{Margin} is collateral posted by one counterparty to the one that collects it, to cover potential losses on
derivatives (slide~22).
\begin{itemize}
  \item \emph{Variation margin} (VM) covers the current present value.
  \item \emph{Initial margin} (IM) covers the potential change of value between the counterparty's default and the
        time $t=\delta t$ at which the surviving party can exit the position. The period $\delta t$ is the
        \emph{margin period of risk} (MPOR).
\end{itemize}
IM is sized as a VaR (or ES) of the value change over the MPOR \ts{24}{10}{38:39}.
\end{definition}

Slide~23 shows it as a picture: a value path up to the default, a fan of possible paths over the MPOR, and a
distribution at close-out whose upper tail is covered by the IM collected and whose lower tail is the IM posted.
IM is held in segregated accounts and cannot be used for trading, so it is a real cost (the MVA of
\cref{ch08:sec-cva}) \ts{24}{10}{1:12:07}.

There are two kinds of IM \ts{24}{10}{38:39}. \emph{Cleared} IM is set by each CCP's own model and has existed for
decades. \emph{Bilateral} IM on uncleared trades was introduced by the post-crisis rules (Basel~III/BCBS--IOSCO) and
phased in from 2016; the industry standard model is ISDA's \emph{Standard Initial Margin Model} (SIMM).

\begin{exercise}[Tang's puzzle with collateral (not from the course)]
Redo the puzzle of \cref{ch08:sec-cva} (value $0\to+100$, default with $R=50\%$) when the trade is under a CSA with
daily VM and zero threshold, the counterparty stops paying VM at default, and the value moves from $100$ to $104$
during a 10-day MPOR. What is the loss with and without an initial margin of $5$? What part of the original $\$50$M
loss does VM remove, and what part does IM remove?
\end{exercise}

\subsection{The SIMM as a delta--gamma VaR}
SIMM is an exposure-based proxy for a 10-business-day ($\delta t=14/365$ in calendar time), $99\%$ VaR
\ts{24}{10}{39:41}. Let $P(t)$ be the portfolio value to the surviving party, the counterparty defaulting at $t=0$,
with $\mathrm{VM}(0)=P(0)$. Ignoring path dependence, cash flows in $(0,\delta t)$ and discounting, and letting only
the market state $\mathcal M$ move by $\Delta\mathcal M$ over the MPOR, a second-order expansion gives
\[
  P(\delta t)\approx P(0)+\bm\delta\T\Delta\mathcal M+\tfrac12\,\Delta\mathcal M\T\gamma\,\Delta\mathcal M ,
\]
with $\bm\delta$ the vector of deltas and $\gamma$ the (symmetric) gamma matrix. The loss is the \emph{increase}
$P(\delta t)-P(0)$, since the collateral is frozen at $P(0)$, and IM should satisfy
$\Prob\bigl(P(\delta t)-P(0)<\mathrm{IM}\bigr)=0.99$.

\begin{proposition}[Moments of the delta--gamma P\&L]\label{ch08:prop-simm}
If $\Delta\mathcal M\sim N(0,C)$, then $\Delta P=\bm\delta\T\Delta\mathcal M+\frac12\Delta\mathcal M\T\gamma\Delta\mathcal M$ has
\[
  \E[\Delta P]=\tfrac12\tr(\gamma C),\qquad \Var(\Delta P)=\bm\delta\T C\bm\delta+\tfrac12\tr\bigl((\gamma C)^2\bigr).
\]
Approximating $\Delta P$ by a normal with these moments,
$\mathrm{IM}\approx\frac12\tr(\gamma C)+2.326\sqrt{\bm\delta\T C\bm\delta+\frac12\tr((\gamma C)^2)}$.
\end{proposition}
\begin{proof}
Write $X=\Delta\mathcal M$. $\E[X\T\gamma X]=\sum_{ij}\gamma_{ij}C_{ij}=\tr(\gamma C)$. By Isserlis' (Wick's) theorem,
$\E[X_iX_jX_kX_l]=C_{ij}C_{kl}+C_{ik}C_{jl}+C_{il}C_{jk}$, so
\[\E[(X\T\gamma X)^2]=(\tr\gamma C)^2+2\tr(\gamma C\gamma C),\qquad \Var(X\T\gamma X)=2\tr((\gamma C)^2);\]
the factor
$\frac14$ gives $\frac12\tr((\gamma C)^2)$. The cross term $\Cov(\bm\delta\T X,X\T\gamma X)$ involves third moments of a
centred Gaussian and vanishes.
\end{proof}
Without gamma this is exactly the model-building VaR of \cref{ch08:sec-estimation} with $T=10$, $\beta=0.99$:
$\mathrm{IM}=2.326\sqrt{10}\sqrt{\bm\delta\T C_1\bm\delta}$ for a daily covariance $C_1$ \ts{24}{10}{39:41}. The
gamma (convexity) terms are the ``two additional terms'' of the slide: SIMM is ``basically VaR''.

\begin{coursenote}[Caveat on slide~25 of \file{lec10}]
The slide gives $\mathrm{IM}\approx2.326\sqrt{10}\sqrt{\bm\delta\T C\bm\delta}+\frac12\sqrt{10}\tr(\gamma C)
+\lambda\sqrt{\frac12\tr((\gamma C)^2)}$. Two remarks.
\begin{itemize}
  \item Splitting the square root of a sum into a delta term and a
      curvature term, with a separately calibrated multiplier $\lambda$, is a modelling choice (SIMM treats delta and
      curvature margins separately), not a consequence of the derivation; moreover $\Delta P$ is not Gaussian, so
      $\lambda$ absorbs the skew of the quadratic term.
  \item The time scaling is inconsistent: if $C$ is the \emph{daily}
      covariance, the 10-day covariance is $10C$, so the mean term is $\frac12\tr(\gamma\,10C)=5\tr(\gamma C)$, not
      $\frac12\sqrt{10}\tr(\gamma C)$, and the last term scales by $10$, not by $1$. Deltas scale like $\sqrt T$, gammas
      like $T$.
\end{itemize}
\end{coursenote}

\begin{exercise}[Delta--gamma margin (not from the course)]
\leavevmode
\begin{enumerate}[(a)]
  \item Prove $\Var(X\T\gamma X)=2\tr((\gamma C)^2)$ for $X\sim N(0,C)$ and symmetric $\gamma$.
  \item One risk factor with
      daily variance $c$, delta $\delta$ and gamma $\gamma$: write mean and variance of the 10-day $\Delta P$ and check how each
      term scales with the horizon.
  \item For a short option position ($\gamma<0$ in the value to the surviving party) discuss
      whether the normal approximation over- or underestimates the $99\%$ quantile.
\end{enumerate}
\end{exercise}

\subsection{SIMM in practice}
The formulation above has too many risk factors (so $C$, $\bm\delta$, $\gamma$ are huge) and $\gamma$ is hard to
compute. SIMM therefore simplifies further (slides 26--27): curvature is proxied from \emph{vega},
$\gamma\propto\mathcal V/(T\sigma_N)$; risk factors are grouped into \emph{buckets} (e.g.\ currencies) and \emph{risk
classes} (IR, FX, equity, commodity, qualifying and non-qualifying credit); sensitivities are delta, vega and
curvature; the results are aggregated up a tree with prescribed correlations, and the \emph{product classes}
(rates/FX, credit, equity, commodity) are simply added with no netting between them. Additional terms handle
thresholds and concentration, and when several regulators apply the posted IM is the maximum over them.

\paragraph{Calibration} \ts{24}{10}{40:43}. The correlations and risk weights, which cover all asset classes, are
recalibrated on historical data. In the lecture (October 2024) SIMM~2.7, due in early December 2024, was expected to
be the first calibration without the 2008 or COVID periods in its window, cutting SIMM IM by roughly $5$--$20\%$
and releasing on the order of $\$40$bn of margin. When the 2020 shocks arrived, the annual calibration made SIMM
underestimate risk for more than a year; ISDA, which owns the model and runs the calibration with the banks, then
moved to semi-annual recalibration \ts{24}{10}{1:13:11}.

\paragraph{Large positions} (slide~28). Big positions take longer to unwind, which is liquidity risk again. One
approach scales the MPOR with the position size, $\Delta T=\Delta T_0\max(1,N/N_0)$; since IM grows like
$N\sqrt{\Delta T}$, beyond $N_0$ one gets $\mathrm{IM}\propto N^{3/2}$. Another adds a concentration add-on
$W(\bm\delta)\,\bm\delta\T C\bm\delta$ with $W$ increasing (e.g.\ piecewise linear). Both keep the IM convex in the
position.

% ==================================================================
\section{Multilateral optimisation of initial margin}\label{ch08:sec-opt}
This is the main topic of the 2024 lecture and the business of Quantile.

\subsection{The problem}
Consider parties $\mathcal P=\{A,B,C,D,E,\dots\}$ (banks) with bilateral portfolios against each other and cleared
portfolios at one or more CCPs, all posting IM (slide~29) \ts{24}{10}{41:48}. The goal is to find a set of new
\emph{hedge trades} between pairs of parties that reduces the total IM in the system (slide~30)
\ts{24}{10}{42:50}. There are no hedge variables with the CCP: one can only book trades between two banks, but if
the chosen instrument is clearable (a vanilla swap) the trade is cleared and changes the cleared IM instead of the
bilateral one \ts{24}{10}{43:55}.

\begin{definition}[Margin minimisation problem, slide 31]\label{ch08:def-mmp}
Let $\mathcal H$ be the set of hedge instruments, $\mathcal M$ the market risk factors, $D^0_{p,q,m}$ the initial
sensitivity of party $p$'s portfolio facing $q$ to factor $m$, and $D_{h,p,q,m}$ the sensitivity of one unit of hedge
$h$. The unknowns are the amounts $x_{h,p,q}$ of hedge $h$ that $p$ trades with $q$. Solve
\begin{align*}
  \min_{x}\quad & \sum_{p\in\mathcal P}\Bigl(\sum_{q\ne p}\mathrm{SIMM}_{p,q}(D_{p,q})
                   +\sum_{c\in\mathcal{CCP}}\mathrm{IM}_{p,c}(D_{p,c})\Bigr)\\
  \text{s.t.}\quad & D_{p,q,m}=D^0_{p,q,m}+\textstyle\sum_{h\in\mathcal H}x_{h,p,q}D_{h,p,q,m} && \forall m\in\mathcal M,\\
  & x_{h,p,q}=-x_{h,q,p} && \forall h,\ p\ne q \quad\text{(symmetry)},\\
  & \textstyle\sum_{q}x_{h,p,q}=0 && \forall h,\ p \quad\text{(flatness)},\\
  & T^-_{p,q,m}\le\textstyle\sum_hx_{h,p,q}D_{h,p,q,m}\le T^+_{p,q,m} && \forall m,\ p\ne q\quad\text{(risk limits)}.
\end{align*}
\end{definition}
The objective is the total margin in the system: bilateral SIMM summed over all pairs, plus cleared IM summed over
CCPs \ts{24}{10}{43:55}. The new risk is the old risk plus the sum of the hedges times their risk
\ts{24}{10}{44:59}. The constraints:
\begin{itemize}
  \item \textbf{Symmetry}: if $p$ sells something to $q$, $q$ must buy exactly that from $p$. This was the
        constraint the banks worried about most: telling Goldman to buy 100 from JPMorgan while telling JPMorgan
        to sell 200 to Goldman would only be discovered when the trade is booked \ts{24}{10}{46:02}.
  \item \textbf{Flatness} (cash-flow neutrality): whatever $p$ sells to one party it buys from others, so each
        party's total position in every instrument, and hence its market risk, is unchanged. Only
        \emph{counterparty} risk is reshuffled, and nobody makes or loses money from the system itself. It is not
        strictly necessary, but relaxing it makes everything much harder \ts{24}{10}{47:02}.
  \item \textbf{Risk limits}: each party can bound how much the risk to each factor against each counterparty may
        change \ts{24}{10}{48:06}.
\end{itemize}

\begin{coursenote}[Notation on slide~31]
The slide writes $\mathrm{IM}_{p,ccp}(D_{p,q})$; the argument should be the cleared sensitivity $D_{p,ccp}$. The
risk update is written for bilateral risk only; for clearable hedge types $h\in\mathcal H_c$ the update goes to the
CCP, $D_{p,c,m}=D^0_{p,c,m}+\sum_{h\in\mathcal H_c}\sum_qx_{h,p,q}D_{h,m}$, which is how cleared IM enters the
optimisation \ts{24}{10}{42:50}.
\end{coursenote}

\subsection{Toy model: the median rule}
\begin{figure}[ht]
\centering
\begin{tikzpicture}[>=Stealth,every node/.style={font=\small},
  pt/.style={circle,draw,thick,minimum size=6mm,inner sep=0pt},
  lab/.style={fill=white,inner sep=1pt,font=\footnotesize}]
\foreach \dx/\ab/\bc/\ca/\ttl in {0/{-2}/{-9}/{-10}/{initial (margin 21)},
                                  4.2/{+9}/{+9}/{+9}/{hedges},
                                  8.4/{7}/{0}/{-1}/{result (margin 8)}}{
  \begin{scope}[xshift=\dx cm]
    \node[pt] (A) at (1.2,2.0) {A}; \node[pt] (B) at (0,0) {B}; \node[pt] (C) at (2.4,0) {C};
    \draw[->,thick,defcol] (A) -- node[lab]{$\ab$} (B);
    \draw[->,thick,defcol] (B) -- node[lab]{$\bc$} (C);
    \draw[->,thick,defcol] (C) -- node[lab]{$\ca$} (A);
    \node at (1.2,-0.7) {\ttl};
  \end{scope}}
\node at (3.3,0.9) {\Large $+$};
\node at (7.5,0.9) {\Large $=$};
\end{tikzpicture}
\caption{The toy model of slide~32. An arrow $p\to q$ labelled $x$ means that $p$'s position against $q$ is $x$
(and $q$'s against $p$ is $-x$). Moving $t=9$ units around the cycle $A\to B\to C\to A$ satisfies symmetry and
flatness by construction and reduces the total $|x|$ from $2+9+10=21$ to $7+0+1=8$.}\label{ch08:fig-triangle}
\end{figure}

\begin{proposition}[Median rule for a triangle]\label{ch08:prop-median}
Three parties, one risk factor, one hedge instrument, bilateral margin proportional to $|$net position$|$ (the
one-factor SIMM $2.326\sqrt{10}\,\sigma|D|$). Let $a_1,a_2,a_3$ be the initial positions oriented around the cycle
$A\to B\to C\to A$. Then every hedge satisfying symmetry and flatness adds the same amount $t$ to all three oriented
edges, and the total margin $\sum_k|a_k+t|$ is minimised at $t=-\operatorname{median}(a_1,a_2,a_3)$.
\end{proposition}
\begin{proof}
Flatness at A gives $x_{AB}+x_{AC}=0$, i.e.\ $x_{AB}=x_{CA}$ by symmetry; at B, $x_{BA}+x_{BC}=0$, i.e.\
$x_{BC}=x_{AB}$. So $x_{AB}=x_{BC}=x_{CA}=t$. The function $g(t)=\sum_k|a_k+t|$ is convex and piecewise linear with
slope $\#\{k:a_k+t>0\}-\#\{k:a_k+t<0\}$, which changes sign at $t=-\operatorname{median}(a)$.
\end{proof}
In the example of slide~32, $a=(-2,-9,-10)$, the median is $-9$, and moving $9$ units around the triangle turns the
positions into $(7,0,-1)$, cutting the margin from $21$ to $8$ \ts{24}{10}{48:06}. Any network can be decomposed into
triangles, and one can think of the optimisation as moving the median around each triangle. Since edges are shared
between triangles and there are additional constraints, this is only a heuristic and not the real solution method
\ts{24}{10}{49:07}.

\begin{coursenote}[Erratum: result panel of slide~32]
In the ``result'' triangle the labels on the A--B edge are $-7$ next to A and $7$ next to B. Adding the ``initial''
($-2$ at A, $2$ at B) and ``hedges'' ($9$ at A, $-9$ at B) panels gives $+7$ at A and $-7$ at B. The margin, which
depends only on $|x|$, is unaffected.
\end{coursenote}

\begin{intuition}[Hedges are circulations on a graph]
Think of $x_{h,\cdot,\cdot}$ as a \emph{flow} on the complete graph of parties: symmetry says it is an antisymmetric
edge function, and flatness says its divergence vanishes at every node (Kirchhoff's current law). Divergence-free
flows are circulations, i.e.\ the cycle space of the graph, of dimension $E-V+1$; for $n$ fully connected parties
this is $\binom n2-n+1=\binom{n-1}{2}$, spanned by triangles. For a triangle it is one number, the $t$ above. The
optimisation moves risk around loops without creating or destroying it at any node.
\end{intuition}

\begin{exercise}[Toy margin optimisation (not from the course)]
\leavevmode
\begin{enumerate}[(a)]
  \item Four parties on a cycle $A\to B\to C\to D\to A$ with oriented positions $a=(3,-5,8,-1)$: find the optimal
      circulation $t$ and the margin saving.
  \item Now add the chord A--C with position $2$: parametrise all hedges satisfying
      symmetry and flatness (the cycle space has dimension $2$) and minimise the total $\sum|x|$.
  \item For the fairness table
      of \cref{ch08:sec-opt}, check the least-squares criterion and propose a different fairness criterion (e.g.\ max--min
      of relative savings) that gives a different answer on some example.
\end{enumerate}
\end{exercise}

\subsection{Solving it in practice}
In reality one writes down the problem and hands it to a numerical solver \ts{24}{10}{50:07}. Some lessons:
\begin{itemize}
  \item \textbf{Convexity is everything.} Minimising ES is piecewise linear, SIMM is \emph{nearly} convex and
        concentration add-ons are \emph{potentially} convex (slide~33). Using general non-linear solvers did not
        work well; it is much easier to approximate the objective by a genuinely convex function and use a mature
        convex solver. Do not write your own \ts{24}{10}{51:07}. How hard a problem is depends on its structure
        and sparsity much more than on its size.
  \item \textbf{Modelling languages} separate the statement of a problem from the numerics. Quantile uses Gurobi;
        the slide's example is the integer program $\max 5x_1+8x_2$ s.t.\ $x_1+x_2\le5$, $3x_1+7x_2\le25$,
        $x_i\in\N$, solved in 10\,ms with optimum $31$ at $(3,2)$ \ts{24}{10}{52:08}.
  \item \textbf{Size.} The FX-delta problem (slide~36): 25 parties, 30 currencies, cleared and uncleared
        non-deliverable forwards (NDFs), one short maturity, about $18{,}000$ variables ($25\cdot24$ ordered pairs
        $\times30$ currencies) and $40{,}000$ constraints (pairwise risk limits, pairwise IM-increase limits, notional
        efficiency, symmetry, flatness, lot sizes), with SIMM plus cleared IM as a convex objective
        \ts{24}{10}{53:14}. NDFs are chosen because each has exposure to exactly one risk factor, one currency pair.
  \item \textbf{Linear algebra.} Interior-point methods take Newton steps: they find the zero of the gradient, and
        each step requires solving with the Hessian, naively $O(n^3)\approx10^{13}$ flops for $n\sim40{,}000$
        \ts{24}{10}{54:18}. In practice one does far better: the matrix is positive semidefinite (Cholesky), very
        sparse (well under $0.1\%$ non-zero) and structured. A few global rows (symmetry) and party rows (flatness)
        sit on top of a block-diagonal part, the $36{,}000$ risk constraints of individual $(p,q)$ pairs, which is
        purely diagonal when each hedge has one risk factor. This \emph{Dantzig--Wolfe} (block-angular) structure
        is exploited by the solvers \ts{24}{10}{55:18}.
  \item \textbf{Presolve} removes redundant rows ($x_1\ge1$ is implied by $x_1\ge2$) and here halved the problem;
        the barrier method then converged in 58 iterations and 11\,s, thanks to the quadratic convergence of Newton's
        method \ts{24}{10}{59:29}.
\end{itemize}

\paragraph{Nearly parallel risks.} The real difficulties are numerical \ts{24}{10}{56:22}. Two trades with almost
identical risk, e.g.\ exposures $3E$ and $E$ to two risk factors, give nearly singular systems (slide~39):
\[
  \begin{pmatrix}1&1\\ \tfrac13&\tfrac13\end{pmatrix}\binom{x_1}{y_1}=\binom{1}{\tfrac13}\ \text{has the solutions } y_1=1-x_1,
  \quad\text{but}\quad
  \begin{pmatrix}1&1\\ 0.3333&0.333\end{pmatrix}\binom{x_1}{y_1}=\binom{1}{b}
\]
has the unique solution $(0,1)$ for $b=0.333$, $(1,0)$ for $b=0.3333$, and $(-110,111)$ for $b=0.3$
\ts{24}{10}{57:25}. The determinant is $-3\times10^{-4}$ and the condition number about $7\times10^3$, so rounding in
the fourth digit moves the solution by $O(1)$ or even by hundreds. With a sign constraint $x_i\ge0$ the last system
is declared \emph{infeasible} purely because of numerical noise \ts{24}{10}{58:28}. Most of the practical work is
detecting and removing such situations, not solving convex problems in general.

\paragraph{Results.} On a 5-party and a 20-party system with one risk factor and no constraints, almost all margin
disappears, and every party ends up either long or short the factor against all its counterparties, never both:
all risk has been netted (slides 41--42) \ts{24}{10}{1:00:32}. With real constraints (slide~43) the picture is less
clean. The empty diagonal is ``margin against oneself'', and a block in the middle turned out to be the affiliates
of one large bank, which do not exchange margin with each other \ts{24}{10}{1:01:35}. A second application
(slide~44) is interest-rate margin: 15 parties, USD/EUR/GBP, swaptions and cleared swaps, about $5000$ variables
and $10{,}000$ constraints, where each trade loads on several risk factors (a swap expiring in 12 years is exposed to
the 10Y and 20Y swap rates).

\begin{exercise}[Rounding and integrality (not from the course)]
For the Gurobi example, $\max 5x_1+8x_2$ subject to $x_1+x_2\le5$, $3x_1+7x_2\le25$, $x_1,x_2\ge0$:
\begin{enumerate}[(a)]
  \item solve the LP
      relaxation (optimum $32.5$ at $(2.5,2.5)$);
  \item show that rounding the LP solution gives either an infeasible point or a
      poor one;
  \item find the integer optimum $31$ by enumeration. Relate this to the rounding of notionals in
      \cref{ch08:sec-opt}.
\end{enumerate}
\end{exercise}

\begin{exercise}[Nearly parallel risks (not from the course)]
For $A_\varepsilon=\bigl(\begin{smallmatrix}1&1\\ 1/3+\varepsilon&1/3\end{smallmatrix}\bigr)$ solve $A_\varepsilon(x,y)\T=(1,b)\T$
and show that $x=(b-1/3)/\varepsilon$. Compute the condition number of $A_\varepsilon$ for small $\varepsilon$ and explain the
three solutions of slide~39. How would you regularise the problem (e.g.\ a small penalty $\eta\|x\|^2$)?
\end{exercise}

\subsection{Nice and fair solutions}
A solution must be \emph{feasible} (all constraints hold) and \emph{optimal} (a good IM saving); in practice it
must also be \emph{nice} and \emph{fair} \ts{24}{10}{1:02:40}.

\paragraph{Nice.} Traders want round notionals (e.g.\ multiples of a million), few trades (\emph{notional
efficiency}), trade packages, and possibly more trades moved into CCPs for reasons beyond IM. The worst headache
came from NDFs: a USD notional and a EUR notional define the FX rate as their ratio, to 15 digits, but banks' booking
systems keep only 6--8 decimals, so the notionals had to be chosen as round dollar amounts that give exact euro
amounts at a 4- or 6-decimal rate \ts{24}{10}{1:03:41}. Rounding cannot be done afterwards. If a party sells 8 to one
counterparty and buys 4 from each of two others, rounding to the nearest 10 gives $10$, $0$, $0$ and breaks
flatness. Rounding must be part of the optimisation, which makes it a \emph{mixed-integer} problem, much harder:
the FX rounding step of slide~46 had about $10^4$ integer variables and took about a minute of branch-and-bound
\ts{24}{10}{1:04:43}.

\paragraph{Fair.} What does fairness mean for an optimiser \ts{24}{10}{1:05:47}? In the four-party example of
slide~47 three solutions give the same total saving:
\begin{center}\small
\begin{tabular}{@{}lcccc@{}}
\toprule
party & scenario 1 & scenario 2 & scenario 3 & single-party best\\
\midrule
A & 4 & 0 & 2 & 4\\
B & 4 & 4 & 4 & 4\\
C & 4 & 4 & 4 & 4\\
D & 0 & 4 & 2 & 4\\
\bottomrule
\end{tabular}
\end{center}
Scenarios 1 and 2 leave one party with nothing; scenario 3 splits the saving, but one could still ask why B and C get
their full 4 \ts{24}{10}{1:06:47}. One proposal is a two-step algorithm. First optimise for each party alone, keeping
all constraints, to get the best achievable $\mathrm{IM}_p^{\text{best}}$; then solve
$\min\sum_p(\mathrm{IM}_p-\mathrm{IM}_p^{\text{best}})^2$ \ts{24}{10}{1:07:53}. Here the squared shortfalls are
$16$, $16$ and $4+4=8$, so the algorithm picks scenario 3 \ts{24}{10}{1:09:00}. Fairness has no established definition
yet; Shepherd considers it the most interesting open problem in this area.

\paragraph{Network effect.} Across real runs, both the saving and the saving as a fraction of the initial IM grow with
the number of participating parties (slide~49): bigger networks benefit everyone \ts{24}{10}{1:10:04}.

\begin{finance}[Optimisation versus compression (not discussed in class)]
\emph{Portfolio compression} terminates redundant offsetting trades (e.g.\ circles of swaps A--B--C--A) and replaces
them with fewer trades of smaller gross notional and the same net risk. The IM optimisation of this lecture instead
\emph{adds} new trades that redistribute risk between pairs of counterparties. Both keep every party's market risk
unchanged (the flatness constraint) and exploit the same cycle structure of the network.
\end{finance}

% ==================================================================
\section{Enterprise-level modelling and martingale tools}\label{ch08:sec-martingale}
The second half of the 2013 lecture turns to the models behind CVA.

\subsection{Trade-level versus enterprise-level models}
A \emph{trade-level} model prices each trade independently (PV and Greeks), and the portfolio PV and Greeks are the
sums: this is what ``derivatives model'' usually means \ts{13}{26}{44:53}. \emph{Enterprise-level} models handle the
quantities that are \emph{not} linear in the trades: CVA, funding, capital (risk-weighted assets), liquidity
\ts{13}{26}{45:57}. They must simulate all trades and all market factors of a portfolio consistently, with a proper
joint distribution, while reusing the trade-level models to value each trade at future dates.

There is also feedback to the trade-level model \ts{13}{26}{47:02}. A cancellable swap facing a counterparty close to
default has an exercise boundary that ignores counterparty credit; cancelling early reduces or removes the CVA. In
practice the CVA desk pays the trading desk the CVA saving, which the trade model can take as an input to its
exercise decision \ts{13}{26}{48:05}.

Putting models together is numerically delicate: some trades are priced by PDE, some by simulation, and a joint
simulation introduces errors \ts{13}{26}{50:14}. Tang's remedy is built on martingales: the arbitrage-free price of
every traded quantity is a martingale under a suitable measure, and those conditions can be tested and enforced.

\subsection{Martingale measures and change of numeraire}
By Harrison and Pliska, if $Y_t$ is the price of a traded asset with no intermediate cash flows and $N_t>0$ is a
\emph{numeraire} (another such price), then $Y_t/N_t$ is a martingale under the measure $\Prob_N$ associated with
$N$ \ts{13}{26}{52:18}. The standard pairs (slides 13--15; see \cref{ch:bs,ch:hjm,ch:quanto}) are:
\begin{center}\small
\begin{tabular}{@{}lll@{}}
\toprule
numeraire & measure & martingale\\
\midrule
$\beta(t)=e^{\int_0^tr_u\dd u}$ & risk-neutral $\Q$ & $Y_t/\beta(t)$\\
$P(t,T)$ (zero-coupon bond) & $T$-forward $\Prob_T$ & forward price $F_Y(t,T)=Y_t/P(t,T)$\\
$P(t,T_{i+1})$ & $\Prob_{T_{i+1}}$ & forward LIBOR $L_i(t)=\frac{1}{\delta_i}\bigl(\frac{P(t,T_i)}{P(t,T_{i+1})}-1\bigr)$\\
$A_{i,N}(t)=\sum_{j=i+1}^N\delta_{j-1}P(t,T_j)$ & annuity $\Prob_{A_{i,N}}$ & swap rate $S_{i,N}(t)=\frac{P(t,T_i)-P(t,T_N)}{A_{i,N}(t)}$\\
$P^D(t,T)$ (domestic) & $\Prob_{T,D}$ & FX forward $F_{FX}(t,T)=\frac{S_{FX}(t)P^F(t,T)}{P^D(t,T)}$\\
\bottomrule
\end{tabular}
\end{center}
The pattern is always the same: write the quantity as a ratio of two traded assets \ts{13}{26}{55:28}. LIBOR is (a
multiple of) a portfolio of two bonds divided by the bond $P(t,T_{i+1})$ \ts{13}{26}{53:23}; the swap rate is a
portfolio of two bonds divided by the annuity \ts{13}{26}{54:27}; the FX forward is interest-rate parity, the foreign
bond converted into domestic currency divided by the domestic bond.

\begin{proposition}[Change of numeraire, slide 16]\label{ch08:prop-numeraire}
Let $N_1,N_2>0$ be numeraires and $Y_T$ an $\mathcal F_T$-measurable payoff. Then
\[
  \E_t^{N_1}[Y_T]=\frac{N_2(t)}{N_1(t)}\;\E_t^{N_2}\!\left[Y_T\,\frac{N_1(T)}{N_2(T)}\right],\qquad
  \text{i.e.}\quad \left.\frac{\dd\Prob_{N_1}}{\dd\Prob_{N_2}}\right|_{\mathcal F_T}=\frac{N_1(T)/N_1(0)}{N_2(T)/N_2(0)} .
\]
\end{proposition}
\begin{proof}
The price $X_t$ of the claim paying $X_T=Y_TN_1(T)$ does not depend on the measure used to compute it
\ts{13}{26}{57:32}:
$X_t=N_1(t)\E_t^{N_1}[X_T/N_1(T)]=N_2(t)\E_t^{N_2}[X_T/N_2(T)]$. Substitute $X_T$ and divide by $N_1(t)$.
\end{proof}
With $N_1=P(\cdot,T_i)$, $N_2=P(\cdot,T_{i+1})$ and $N_1(T)/N_2(T)=1+\delta_iL_i(T)$ for $T\le T_i$:
$\E_t^{T_i}[Y_T]=\E_t^{T_{i+1}}\bigl[Y_T(1+\delta_iL_i(T))\bigr]/(1+\delta_iL_i(t))$. This is the change of measure
behind the LIBOR market model \ts{13}{26}{56:29}.

\paragraph{Credit: the survival measure} \ts{13}{26}{58:38}. By analogy with the swap rate one expects the forward
CDS par spread $\tilde S_j(t,T_i,T_{i+1})$ of name $j$ to be a martingale under a ``risky annuity'' measure. The
catch is that the risky annuity $\tilde A_j$ (the value per unit spread of the premium leg, which pays only while
the name survives, i.e.\ a portfolio of zero-recovery defaultable bonds) drops to zero at default, and a numeraire
must stay positive \ts{13}{26}{59:40}. Sch\"onbucher's solution is to ignore the default states
\ts{13}{26}{1:00:45}: since $\ind_{\{t<\tau_j\}}\tilde S_j\tilde A_j$ is the value of the premium leg, a traded asset,
$\ind_{\{t<\tau_j\}}\tilde S_j\tilde A_j/\beta(t)$ is a $\Q$-martingale, and the \emph{survival measure}
\[
  \left.\frac{\dd\Prob_{\tilde A_j}}{\dd\Q}\right|_{t}=\frac{\ind_{\{t<\tau_j\}}\tilde A_j(t,T_i,T_{i+1})}{\beta(t)\,\tilde A_j(0,T_i,T_{i+1})}
\]
makes $\tilde S_j(t,T_i,T_{i+1})$ a martingale \ts{13}{26}{1:03:55}. $\Prob_{\tilde A_j}$ is absolutely continuous
with respect to $\Q$ but \emph{not equivalent} (it gives zero mass to default before $t$). One difficulty is traded
for another, but the second is easier, provided what happens at default is modelled separately
\ts{13}{26}{1:01:49}.

\begin{exercise}[Change of numeraire (not from the course)]
\leavevmode
\begin{enumerate}[(a)]
  \item Prove \cref{ch08:prop-numeraire} and deduce that $L_i(t)$ is a $\Prob_{T_{i+1}}$-martingale.
  \item Show that the swap
      rate $S_{i,N}(t)$ is a martingale under the annuity measure and that it is a weighted average of forward LIBORs with
      weights $\delta_jP(t,T_{j+1})/A_{i,N}(t)$.
  \item Verify \cref{ch08:prop-mginterp} and show by an example that ordinary
      linear interpolation in maturity $T$ is not a martingale in general.
\end{enumerate}
\end{exercise}

\subsection{Martingale testing, resampling and interpolation}
\begin{itemize}
  \item \textbf{Martingale testing}: check in the numerical implementation that known martingale relations such as
        $F(s)=\E_s[F(t)]$ hold; usually they fail slightly \ts{13}{26}{1:04:58}.
  \item \textbf{Martingale resampling}: transform the simulated variables (linearly or log-linearly) so that they
        satisfy the relations \emph{exactly}, keeping a variance metric fixed. The quadratic resampling of slide~19
        is $X=\frac{\mathrm{Std}(X)}{\mathrm{Std}(X_0)}\bigl(X_0-\E[X_0]\bigr)+\E[X]$, with $X_0$ the raw simulated values,
        $\E[X]$ the martingale target and $\mathrm{Std}(X)$ the target dispersion \ts{13}{26}{1:05:59}. This is moment
        matching, familiar from variance reduction in Monte Carlo.
  \item \textbf{Martingale interpolation}: models simulate only a few tenors (say 1Y and 5Y LIBOR). Interpolating a
        missing tenor (3Y) linearly in maturity does \emph{not} preserve the martingale property; interpolating
        linearly using the time-$s$ values as the independent variable does \ts{13}{26}{1:07:01}.
\end{itemize}

\begin{proposition}[Martingale interpolation]\label{ch08:prop-mginterp}
Let $M(\cdot,T_1)$ and $M(\cdot,T_2)$ be martingales, $M(s,T)=\E_s[M(t,T)]$ for $s\le t$, and let the target
$M(s,T_3)$ be known at a reference time $s$ (e.g.\ $s=0$, from today's curve). Define for $t\ge s$
\[
  M(t,T_3)=\frac{M(s,T_3)-M(s,T_2)}{M(s,T_1)-M(s,T_2)}\,M(t,T_1)+\frac{M(s,T_1)-M(s,T_3)}{M(s,T_1)-M(s,T_2)}\,M(t,T_2).
\]
Then $\E_s[M(t,T_3)]=M(s,T_3)$ \ts{13}{26}{1:10:06}.
\end{proposition}
\begin{proof}
The weights are $\mathcal F_s$-measurable and sum to $1$. Taking $\E_s$ replaces $M(t,T_k)$ by $a_k=M(s,T_k)$, giving
$\bigl[(a_3-a_2)a_1+(a_1-a_3)a_2\bigr]/(a_1-a_2)=a_3$ \ts{13}{26}{1:11:10}.
\end{proof}
Equivalently, $M(t,T_3)$ is read off the straight line through the points $(a_1,M(t,T_1))$ and $(a_2,M(t,T_2))$ at the
abscissa $a_3$: linear interpolation with the ``right'' independent variable.

\paragraph{The LIBOR market model} \ts{13}{26}{1:12:13}. Martingale modelling proceeds in three steps: find the
measure under which a quantity is a martingale, represent the martingale as a stochastic integral (log-normal, CEV or
stochastic volatility), then change to a common measure \ts{13}{26}{1:13:17}. For forward LIBORs,
$\dd L_i/L_i=\lambda_i(t)\dd W_t^{T_{i+1}}$ under $\Prob_{T_{i+1}}$. Under the terminal measure $\Prob_{T_N}$ (slide~20),
\[
  \frac{\dd L_i(t)}{L_i(t)}=-\sum_{j=i+1}^{N-1}\frac{\lambda_i\lambda_j\delta_jL_j(t)}{1+\delta_jL_j(t)}\,\rho_{ij}\dd t
  +\lambda_i\dd W_t^{T_N,L_i},\qquad \rho_{ij}=\Corr\bigl(\dd\ln L_i,\dd\ln L_j\bigr),
\]
the full-dimensional model with one Brownian motion per LIBOR and the correlations as inputs \ts{13}{26}{1:15:33}.
The drift follows from \cref{ch08:prop-numeraire} and Girsanov's theorem (\cref{ch:stochcalc}); interest-rate models
of this kind are the subject of \cref{ch:hjm}.

% ==================================================================
\section{Closing remarks of the 2013 course}\label{ch08:sec-closing}

\begin{coursenote}[Tang's conclusion and the end of 18.S096 (2013)]
Tang, who presented from New Jersey because snow had grounded his flight \ts{13}{26}{0:00}, concluded that banks need
enterprise-level models for the non-linear portfolio effects (CVA, funding, liquidity, capital), built on top of
trade-level models with martingale testing, resampling and interpolation. These risks have received much more
attention since the crisis \ts{13}{26}{1:16:34}.

The instructors then closed the course \ts{13}{26}{1:17:37}. A list of past final-paper topics was on the course site:
mostly Black--Scholes, its extensions and manipulations of the Black--Scholes equation, one statistical analysis of
commodity data (recommended as a direction), and numerical and Monte Carlo projects. Students were asked to reflect
on what they had learned and on what they would like to explore next. P.~Kempthorne described the course as
deliberately challenging, meant to give both the mathematical foundations and a real exposure to their applications
through practitioners; he said the final papers he had read so far were good and that the faculty wanted to remain
a resource after the course \ts{13}{26}{1:19:45}. The class would be repeated the following fall in a slightly
different form, students were invited to suggest topics (one value of the course being access to speakers from the
frontier of the industry), and to fill in the course review \ts{13}{26}{1:20:45}.
\end{coursenote}

\begin{coursenote}[2013 versus 2024]
The two lectures reflect a decade of regulation. In 2013 many OTC trades were uncollateralised, and the focus was
pricing and hedging CVA, DVA and funding at portfolio level. By 2024 standard swaps are centrally cleared and
bilateral trades between banks carry daily VM and regulatory IM, so counterparty risk is managed mainly through
collateral: the questions are how much margin (a VaR/ES) and how to reduce it (optimisation), and the cost of
counterparty risk has partly turned into the cost of funding margin. Both speakers stress the same point: the
relevant quantities are non-linear functions of the \emph{whole} portfolio, or even the whole network.
\end{coursenote}

\begin{coursenote}[Reading suggested on the slides]
2024: S.~Boyd and L.~Vandenberghe, \emph{Convex Optimization} (the slide misspells the first name as ``Steven'');
J.~Gregory, \emph{The xVA Challenge}; J.~Hull, \emph{Risk Management and Financial Institutions}; Artzner et al.\
(1999) on coherent risk measures; Rockafellar and Uryasev on optimising CVaR; the BCBS--IOSCO margin framework and
the ISDA SIMM~2.6 methodology; Embrechts and Wang, ``Seven proofs for the subadditivity of expected shortfall'';
``Subadditivity re-examined: the case for Value-at-Risk''. 2013: Harrison--Kreps (1979), Harrison--Pliska (1981),
Jamshidian (1997) on LIBOR and swap market models, Sch\"onbucher (2003) on survival measures, Tang and Li (2007),
Tang and Williams (2010) on funding benefit and cost, Zhu and Pykhtin (2007), ``A guide to modeling counterparty
credit risk''.
\end{coursenote}

% ==================================================================
\begin{keyformulas}
\begin{align*}
&\text{Exposure:} && E(t)=(V(t)-C(t))^+,\\
& && \mathrm{EE}(t)=\E[E(t)],\quad
  \textstyle(\sum_iV_i)^+\le\sum_iV_i^+\\
&\text{CVA:} && \mathrm{CVA}=-\E^{\Q}\bigl[\ind_{\{\tau\le T\}}(1-R)
  (V_{\tau-}-C_{\tau-})^+/\beta_\tau\bigr]
  \approx-\mathrm{EPE}\cdot s\cdot D\\
&\text{DVA:} && \mathrm{DVA}=\E^{\Q}\bigl[\ind_{\{\bar\tau\le T\}}(1-\bar R)
  (C_{\bar\tau-}-V_{\bar\tau-})^+/\beta_{\bar\tau}\bigr]\\
&\text{VaR, ES:} && \mathrm{VaR}_\beta=\inf\{\ell:\Prob(L\le\ell)\ge\beta\},\\
& && \mathrm{ES}_\beta=\tfrac1{1-\beta}\textstyle\int_\beta^1\mathrm{VaR}_u\dd u\\
&\text{Normal:} && \mathrm{VaR}_\beta=\mu+z_\beta\sigma,\\
& && \mathrm{ES}_\beta=\mu+\sigma\varphi(z_\beta)/(1-\beta),\\
& && \sigma_P^2=\bm\delta\T C\bm\delta\\
&\text{Time scaling:} && \text{factor }\sqrt{T+2\textstyle\sum_{k=1}^{T-1}(T-k)\rho^k}\\
& && (\sqrt T\text{ if }\rho=0)\\
&\text{Back-test:} && \Prob(N\ge m)=\textstyle\sum_{k\ge m}\binom nkp^k(1-p)^{n-k}\\
&\text{Rockafellar--Uryasev:} && \mathrm{ES}_\beta=\min_\alpha\bigl\{\alpha+\tfrac{1}{1-\beta}\E[(L-\alpha)^+]\bigr\}\\
&\text{Delta--gamma IM:} && \E[\Delta] P=\tfrac12\tr(\gamma C),\\
& && \Var\Delta P=\bm\delta\T C\bm\delta+\tfrac12\tr((\gamma C)^2)\\
&\text{Margin optim.:} && \min\textstyle\sum_p\bigl(\sum_q\mathrm{SIMM}_{p,q}
  +\sum_c\mathrm{IM}_{p,c}\bigr),\\
& && x_{h,p,q}=-x_{h,q,p},\ \sum_qx_{h,p,q}=0\\
&\text{Change of numeraire:} && \E_t^{N_1}[Y_T]=\tfrac{N_2(t)}{N_1(t)}
  \E_t^{N_2}\bigl[Y_TN_1(T)/N_2(T)\bigr]
\end{align*}
\end{keyformulas}
